\section{統計学的にデータを調べるということ}

  \subsection{運まかせを疑う}

    じゃんけんに負けても, 普通は運が悪かったと思うくらいだろう.
    どの手が出るかは運まかせで, 勝てるのは3回に1回だと思っているからだ.
    ところが, 出した手を何千回も記録してみると, 人の出す手には案外偏りがある.
    ある調査\footnote{%
    芳沢光雄『じゃんけんの算数：樹形図と確率の考えかた』日本評論社, 2006年.
    学生725人・延べ11{,}567回のじゃんけんデータに基づく.
    集計値は長谷実李・青山和裕「じゃんけんの強さとその人の性格特徴の相関性」『日本科学教育学会研究会研究報告』34巻3号, 2019年, 53--58頁, doi:10.14935/jsser.34.3\_53にも再掲されている.}
    によると, グー・チョキ・パーの出た割合はおよそ 35 : 32 : 33 だったそうだ.
    35\%と32\%ならたったの3ポイント差で, ほとんど互角に見える.
    100回くらいのじゃんけんなら, このくらいの差は偶然でもよく出る.
    だが, 11{,}567回も集めた上での3ポイント差となると, これは偶然とは片づけにくい.

    ただ, 実際に勝てるかどうかは別の話だ.
    仮に相手がこの割合のままずっと手を出し続けてくれたとしても, こちらがパーを出し続けて勝てる割合は約35\%.
    相手がまったくでたらめに出すときの約33\%と比べて, 2ポイントほど高いにすぎない.

    さて, ここまでだと「なんだ, たった2ポイントか. もっと勝てる方法はないのか」という感じかもしれない.
    ところが, じゃんけんを何回か続ける状況を考えると, 様子が変わってくる.
    手がかりになるのは, 相手が直前に何を出したかだ.
    でたらめに出しているなら, 1回目と同じ手を2回目にまた出すのは3回に1回のはずだ.
    同じ人が2回続けてじゃんけんをした10{,}833組の記録を見ると, 1回目と同じ手をまた出したのは2{,}465組で, 約23\%しかなかった.
    たとえば, グーであいこになったとしよう.
    相手がこの傾向どおりに手を出すとすれば, 次も続けてグーを出してくるのは23\%くらいしかない.
    それなら, こちらはチョキを出せばいい.
    負けるのは相手がもう一度グーを出したときだけだ.
    当てもなく手を出せば $1/3$ の確率で負けるところを, 23\%ほどまで下げられる.

    だから, 自分はグーを出し続けるのだ, という強固な信念がない限り, あいこになった場合はこの傾向に従ってチョキを出すほうがクレバーだといえる.
    このように, 確実にこうだ, ということはわからないが, こうなる可能性が比較的高いぞ, と言うための根拠を与えてくれるのが統計のよいところの一つだ.
    現実は複雑なのでもごもごした物言いになってしまうが, それが統計の味わいというものである.

  \subsection{確率論で扱うサイコロと統計学で扱うサイコロ}

    確率論と統計学はいずれも平たく言うと「物事がどのくらい起こるのか」を考えるものであり, ともするとどんな違いがあるか, 
    はっきりしないと思われるかもしれない.
    これにはいろいろな表現があると思うが, ここでは簡単に, 確率論はより理論的であり, 統計学はより実践的だ, という違いを挙げておきたい.

    たとえば, 手元にサイコロが一つあって, 暇つぶしに（ちょっと寂しすぎる気もするが）20回振ってみたとしよう.
    1の目が出た回数を数えてみると, なんと10回もあった.
    1の目が出る確率は $1/6$ だから, 20回振れば3回か4回程度になるはずだ.
    半分が1というのは, いくらなんでも多すぎる.
    「このサイコロ, 中に重りでも入っているのではないか」と疑いたくなる.

    とはいえ, まともなサイコロであっても, 20回振って1がきっちり3回や4回出るとは限らない.
    1回しか出ないこともあれば, 6回出ることだってある.
    10回出る確率も, 決してゼロではない.
    極端な話をすれば, 20回すべてが1になる確率さえゼロではないのだ.
    1が10回出るという現象は, 細工のあるサイコロでも, ないサイコロでも起こりうる.
    どれだけサイコロを眺めても, どちらか区別できない.
    すると, 1が10回出たという結果から判断するしかないが, これだけでは細工があるとまでは言い切れない.

    もっとも, 20回すべてが1になるような確率まで気にしていたらきりがないが, やはり10回も出るのは怪しい.
    そう感じるのは, 「まともなサイコロなら, そんなに1ばかり出ないはずだ」と思うからだ.
    それなら, まともなサイコロで1が10回以上出るのがどれくらい珍しいことなのかを計算してみよう（ちょうど暇だ）.
    たとえば1が6回程度なら, 「今日はよく1が出るな」くらいで済ませるだろう.
    実際, まともなサイコロを20回振って1が6回以上出る確率は, 1割ほどある.
    ところが10回以上となると, 確率はおよそ1700分の1にしかならない.
    「たまたま」で片づけるには, さすがに珍しすぎる確率だ.

    ただし, この1700分の1という数値は, 「どの目も $1/6$ の確率で出る」とあらかじめ定めた, まともなサイコロを前提に計算したものだ.
    確率論が扱うのはこのようなサイコロであり, 前提となる条件が初めから決まっているため, 正しく計算すると答えはただ一つに定まる.
    一方, 統計学が扱うのは目の前にある実際のサイコロだ.
    その構造が本当に偏りのないものなのかは分からないし, 何回振ったところで完全に確かめることはできない.
    そこで統計学では, 「絶対に細工がある」と断定する代わりに, 「確率が1700分の1という極めて珍しい出来事が起きたか, さもなければサイコロに細工があるかのどちらかだ」と主張する.
    どちらが真実なのかまでは言及しないが, 1700分の1というおかしな数値を聞けば, 誰もがサイコロの仕込みを疑いたくなるだろう.

    アンケートの回答から社会の意見を推し量ることも, 1万回を超えるじゃんけんの記録から人間の癖を探ることも, このサイコロの例と同じだ.
    確率論は, あらかじめまともな作りだと決めたサイコロを出発点として, どの目がどれくらいの割合で出るかといったことを考える.
    反対に統計学は, 実際に目の前にあるサイコロを振って出た目を出発点として, そのサイコロの構造について何が言えるかを探るのだ.
    出た目からサイコロの構造へさかのぼる過程には, どうしても確実と言い切れない部分が残ってしまう.
    このように統計学は, 不確かな部分をあえて言い切らずに残したまま, 目の前のデータがどれほど珍しいのかを確率というグラデーションで表現する.

  \subsection{このパートで学ぶこと}

    このパートでは, 具体的な例を一つずつたどりながら, 珍しさを表す確率がどこから出てくるのか, それをもとに何がどこまで言えるのかを考えていく.
    そうしていくうちに, もごもごした物言いにも, ちゃんとした理屈があることが見えてくるはずだ.
    完全に納得とはいかないまでも, こういう考え方もありだな, と受け入れてもらえれば十分である.

    まずは, データを平均のような一つの数にまとめるところから始める.
    何が珍しいのかを考えようにも, 何の数字について考えるのかが決まっていなければ話にならないからだ.
    手元のデータについて知りたいだけなら, 平均を計算すればそれで終わりだ.
    ただ, 知りたいのはたいてい, 手元のデータよりもっと大きな全体のほうだ.
    アンケートなら, 知りたいのは回答した人たちの意見ではなく, 社会全体の意見だろう.
    ところが, 手元のデータで計算した値は, 当然データを取り直すと別の値になってしまう.
    どのくらい変わるのかがわからなければ, 手元の平均から全体についてどこまで言っていいのかもわからない.
    それでも, 取り直したときに平均がどのくらいばらつくのかは, 手元のデータだけから計算できてしまう.
    このばらつきをもとに, 「全体の平均はだいたいこのあたりだろう」「二つのグループの差はこれくらい珍しい」といったことを確率で表していく.
    もっとも, 差が珍しいとわかっても, なぜその差が生まれたのかまではわからない.
    原因を確かめるには, 計算よりも, どう比較するかのほうが大事になってくる.
    それに, データの集め方が偏っていたら, せっかくの計算も当てにならない.
    そこでこのパートの最後では, データがどうやって集められたのかまでさかのぼって考える.

    推定や検定, 相関といった, 統計学の入門でかならず出てくる話題は, このパートでもひととおり扱う.
    だから, 読み終えたあとで統計学をもう少し深く学ぼうと思ったときにも, ここで扱った話題が足場になるはずだ.
